
\paragraph{Statement A.4.3 (Proved uniform spectral coefficient estimate).}
Let $M$ be a fixed compact hyperbolic surface, let $l_0,l_1$ be fixed
closed geodesics of lengths $L_0,L_1$, and let $u_j$ be an orthonormal
basis of scalar Laplace eigenfunctions with
$\lambda_j=1/4+t_j^2$, $t_j\geq0$. Set $s_j=1/2+it_j$ and
\[
 P_{j,\ell}=L_0^{-1/2}\int_{l_0}u_j(u)e^{2\pi i\ell u/L_0}\,du,
 \quad P_j'=L_1^{-1/2}\int_{l_1}u_j\,du,\quad
 \kappa_\ell=2\pi\ell/L_0,
\]
\[
 D_s(\kappa)=
 \frac{\Gamma((s+i\kappa)/2)\Gamma((s-i\kappa)/2)}{\Gamma(s)},
 \qquad
 Q_{j,H}=\sum_{|\ell|\leq H}b_\ell
             \frac{D_{s_j}(0)}{D_{s_j}(\kappa_\ell)}P_{j,\ell},
\]
where $|b_\ell|\leq B/(1+|\ell|)$ and $H\geq1$. Then, for $T\geq1$,
\begin{equation}\label{A:uniform-arc-period}
 \sum_{T\leq t_j<2T}|Q_{j,H}|^2
       \ll B^2T\log^2(2H),\qquad
 \sum_{T\leq t_j<2T}|P_j'Q_{j,H}|
       \ll BT\log(2H).
\end{equation}
The constants depend on the surface and geodesics, but not on $T,H$
or the displayed coefficients. Thus the angular sum itself can be
controlled without a power of $H$ in this normalized spectral expression.
This statement does not by itself assert that it is the complete
spectral expression for the sharp count with a uniform remainder.
\begin{proof}
We first record an elementary broad-band restriction estimate, uniformly
in the real frequency $\kappa$:
\begin{equation}\label{A:broad-band-period}
 \sum_{t_j\leq2T}\left|\int_{l_0}u_j(u)e^{i\kappa u}\,du\right|^2
       \ll T.
\end{equation}
Indeed, in finitely many Fermi charts along the parametrized geodesic,
the one-dimensional fundamental theorem of calculus and Cauchy--Schwarz
give, for $0<\delta\ll1$,
\[
 \|f|_{l_0}\|_2^2\ll
 \delta^{-1}\|f\|_{L^2(M)}^2+
 \delta\|\nabla f\|_{L^2(M)}^2.
\]
The finite number of charts accounts for any self-intersections.
For a linear combination of eigenfunctions with $t_j\leq2T$,
$\|\nabla f\|_2^2\leq(4T^2+1/4)\|f\|_2^2$.
Take $\delta=T^{-1}$ (adjust a fixed constant for small $T$), and
apply the adjoint restriction operator to $e^{i\kappa u}$, whose
$L^2(l_0)$ norm is fixed. This proves \eqref{A:broad-band-period}.
The same estimate applies to $P_j'$.
We also use the elementary bound $\#\{t_j\leq2T\}\ll T^2$.
One proof partitions a fixed finite collection of coordinate patches
into $O(T^2)$ cells of diameter $c/T$ and uniformly bounded shape.
If the spectral space had larger dimension, a nonzero member would
have integral zero on every cell. The cell Poincar\'e inequalities would
give $\|f\|_2^2\leq C(c/T)^2\|\nabla f\|_2^2$, contradicting the
spectral bound when the fixed $c$ is sufficiently small.

Reznikov's mean-value estimate
\cite[equation (1.13), proved in Sections 2.3 and 3.1--3.2]{Reznikov2015}
states that the normalized hyperbolic coefficients $a_{j,\ell}$ satisfy
\begin{equation}\label{A:Reznikov-mean}
 \sum_{|\ell|\leq N}|a_{j,\ell}|^2\ll\max(N,1+t_j),
 \qquad P_{j,\ell}=a_{j,\ell}D_{s_j}(\kappa_\ell).
\end{equation}
A fixed scalar from the model functional in his equation (2.3) has
been absorbed into $a_{j,\ell}$. The hypotheses are precisely a fixed
closed geodesic on a compact hyperbolic surface and an $L^2$-normalized
principal-series eigenfunction. His argument includes both components
of the diagonal group; retaining the component that gives scalar
periods only decreases its nonnegative sum. The assertion is uniform
in the eigenfunction as well as in $N$.

Split the angular modes at $|\kappa_\ell|=T/2$.
For $t_j\in[T,2T]$ and the lower modes, Stirling's formula in the
fixed vertical strip $\operatorname{Re}(s_j/2)=1/4$ gives
\[
 \left|\frac{D_{s_j}(0)}{D_{s_j}(\kappa_\ell)}\right|\ll1.
\]
For bounded $T$ this follows from continuity and the absence of gamma
zeros or poles; for large $T$ both $t_j\pm\kappa_\ell$ are comparable
to $T$ and the exponential factors cancel exactly.
Minkowski's inequality in the $j$ variable and
\eqref{A:broad-band-period} bound the square root of the lower-mode
sum by
\[
 \sum_{|\kappa_\ell|\leq T/2}|b_\ell|
 \left(\sum_{T\leq t_j<2T}|P_{j,\ell}|^2\right)^{1/2}
       \ll B\sqrt T\log(2H).
\]

For the higher modes divide the integers into dyadic blocks
$N\leq|\ell|<2N$, with $N\gg T$ up to a fixed geodesic-length factor.
On such a block \eqref{A:Reznikov-mean} and Cauchy--Schwarz give
\[
 \left|\sum_{N\leq|\ell|<2N}b_\ell a_{j,\ell}\right|
 \leq\left(\sum|a_{j,\ell}|^2\right)^{1/2}
       \left(\sum|b_\ell|^2\right)^{1/2}\ll B.
\]
Stirling's formula also gives $|D_{s_j}(0)|\ll T^{-1/2}$.
Since the gamma factors cancel exactly in the definition of $Q$,
the higher-mode contribution to each $Q_{j,H}$ is at most
$BT^{-1/2}\log(2H)$. There are $O(T^2)$ eigenvalues in the band,
so its squared sum is $O(B^2T\log^2(2H))$.
Combining the two ranges proves the first estimate in
\eqref{A:uniform-arc-period}; Cauchy--Schwarz with
$\sum|P_j'|^2\ll T$ proves the second.
\end{proof}

The estimate applies after torsion removal to the fixed cover used in
Statement A.3.1. The original angular frequencies are a fixed-index
subsequence of the cover's frequencies, and the constants in
\eqref{A:Reznikov-mean} change only by that fixed index. This removes
one concrete obstacle in a possible twisted relative trace proof.
A uniform radial transform and an error estimate for its inversion
remain to be established; we do not replace that step by the untwisted
transform of Lekkas--Petridis.
